\needspace{8\baselineskip}
\section{模型假设}

题面已给出的事实（初始温度 $28$ °C、初始水分浓度 $2.55$ kg/kg、阈值 $0.15$ kg/kg 等）
不重复包装为假设。下面六条是题面未完全规定、而建模必须补足的现实简化，
每条都写明它在何处被使用。

（1）轴向与周向输运可忽略，问题退化为一维径向轴对称。侧面积是两端面积的 $12.5$ 倍、
端面只占总换热面积的 $7.4\%$，且题面要求的输出只以“到药材中心的距离”索引，
保留轴向坐标会使同一半径不同高度的结果不同、制表失去唯一解\upcite{adrover2019twod}。

（2）水分迁移由液相扩散主导，通量按 Fick 定律并用附录给出的扩散系数闭合。
附录 2、3、4 直接给出“水分浓度扩散系数”的经验式，等于题面已指定本构形式；
另设毛细或蒸发前沿机制需要题面未提供的相变焓，属于替换题设\upcite{luikov1966capillary,onwude2016pumpkin}。

（3）不显式建模汽化潜热引起的表面降温。附录未给出汽化潜热，代入该量必须自拟数值；
这是本文最主要的物理简化，其后果在第~\ref{sec:validation}~节由对流换热系数的
低敏感性界定\upcite{dedic2002wood}。

（4）药材骨架做仿射收缩，即每个物质点的 $r/R(t)$ 保持不变，该假设仅在问题四启用。
附件 2 只给出外半径一条曲线、没有内部位移信息，仿射收缩是在这一观测条件下
最简且能保证干物质守恒的闭合假设\upcite{adrover2019general,krokida1998shrinkage}；
它并非唯一可能的运动学，若内部收缩不均匀（如表层致密化），
物质点的径向分布会偏离本假设，本文无从由单一半径曲线分辨这种偏离。

（5）对流换热系数与对流传质系数在四问中统一沿用附录 2 的取值。附录 3 与附录 4 只重新给出
密度、比热容、热传导系数与扩散系数，而这两个系数刻画的是烘房空气侧的边界层，
不随药材物性闭合而改变。

（6）$t>14400$ s 的烘房状态由附件 1 末段平台值常数延拓。附件 1 末 $24$ 个点的温度均值为
$50.0095$ °C、含湿量均值为 $0.049986$ kg/kg，波动已收在 $\pm0.26$ °C 内，
用该均值延拓正对应题面“恒温干燥”的字面含义。
